\subsection{Expected Values}
\hrulefill
\vspace{5mm}

The \textbf{expected value} is a weighted average of the possible values of a random variable X with PMF p(x).
It is given by:
\[
    E(X) = \mu_X = \sum_x x p(x)
\]
where the sum is over all possible values of X.

\paragraph*{Names for Expected Value}
\begin{itemize}
    \item Mean
    \item Expected Value
    \item Expectation
    \item $\mu$
    \item $E(X)$
    \item (Bonus) The "center of mass" of the histogram of the PMF.
\end{itemize}

\paragraph*{Example:}
Craps is a dice game in which the players make wagers on the sum of a roll of a pair of dice. One of the wagers 
available to the player is the \textit{field bet}. The field consists of \{5, 6, 7, 8\}. If the sum of the dice is one of these
numbers, you lose your bet. Otherwise you get your bet back plus a payout according to the following table:

\begin{table}[h!]
    \centering
    \begin{tabular}{c|c|c|c|c}
        Winnings & -1 & 1 & 2 & 3 \\
        \hline
        Outcome & field & {3, 4, 9, 10, 11} & 2 & 12 \\
    \end{tabular}
\end{table}

Is the bet favorable to the player?

\paragraph*{Solution:}
Let X be the random variable that gives your winnings on a field bet. The PMF is given by:

\[
\begin{array}{c|c|c|c|c}
    x & -1 & 1 & 2 & 3 \\
    \hline
    p(x) & \frac{16}{36} & \frac{20}{36} & \frac{1}{36} & \frac{1}{36} \\
\end{array}
\]

\begin{align*}
    E(X) &= \sum_x x p(x) \\
    &= -1 \cdot \frac{20}{36} + 1 \cdot \frac{14}{36} + 2 \cdot \frac{1}{36} + 3 \cdot \frac{1}{36} \\
    &= -\frac{20}{36} + \frac{14}{36} + \frac{2}{36} + \frac{3}{36} \\
    &= -\frac{1}{36} \\
    &\approx -0.028
\end{align*}
The expected value is negative, so the bet is unfavorable to the player.

\subsubsection*{Expectation of a Function of a Random Variable}
\vspace{5mm}

If $X$ is a discrete random variable with PMF $p(X)$ and $h(x)$ is any real valued function of X, then the expectation of 
$h(X)$ is given by:
\[
    E[h(X)] = \sum_x h(x) p(x)
\]


\paragraph*{Example:}
Find $E(X^2)$ if X is the sum of two dice.

\[
    E(X^2) = \sum x^2 p(x) 
\]

\pagebreak

Note that this is not always reversible. $E[h(X)] \neq h(E[X])$ except when $h(x)$ is a linear function. Then 
\[
    E(aX + b) = aE[X] + b
\]

\paragraph*{Proof:}
\begin{align*}
    E(aX + b) &= \sum_x (ax + b)p(X) \\
    &= \sum_x ax \cdot p(x) + b \cdot p(x) \\
    &= a \sum_x x p(x) + b \sum_x p(x) \\
    &= aE[X] + b \cdot 1 \\
    &= aE[X] + b
\end{align*}

\paragraph*{Special Cases and Rules:}
\begin{align*}
    a = 0 &: E(b) = b \\
    b = 0 &: E(aX) = aE[X] \\
    E(X + Y) &= E(X) + E(Y) \\
    E(3X - 2Y + b) &= 3E(X) - 2E(Y) + b
\end{align*}

\subsubsection*{Variance of a Discrete Random Variable}
\vspace{5mm}
The \textbf{variance} of a discrete random variable is the "spread" of the random variable around it's mean.
Suppose X is a discrete random variable with PMF p(x) and mean $\mu$. The variance of X, denoted by $\sigma^2$ is given by:
\[
    \sigma^2 =Var(X) = E[(X - \mu)^2] = E[X^2] - \mu^2
\]
Recall here that $\mu = E(X)$. The variance is the mean squared distance of the values of X from their mean.

\pagebreak

\paragraph*{Names for Variance}
\begin{itemize}
    \item Variance
    \item $Var(X)$
    \item $V(X)$
    \item $\sigma^2$
\end{itemize}

\paragraph*{Computation Formula}
\[
    V(X) = E[X^2] - (E[X])^2
\]

\subsubsection*{Standard Deviation}
\vspace{5mm}
The \textbf{standard deviation} of a discrete random variable X is the positive square root of the variance of X. It is denoted by $\sigma$.
\[
    \sigma = SD(X) = \sqrt{Var(X)} = \sqrt{E[X^2] - (E[X])^2}
\]

\paragraph{Note:}
Variance is NOT a linear operator. Let X be a random variable with constants a and b. Then:
\[
    V(aX + b) = a^2 V(X) \implies \sigma_{aX + b} = |a| \sigma_X
\]

\paragraph*{Example:}
Show that the variance of the Bernoulli Random Variable is $p(1-p)$:
\begin{align*}
    V(X) &= E(X^2) - [E(X)]^2 \\
    &= \sum_x x^2 p(x) - \left(\sum_x x p(x)\right)^2 \\
    &= 0^2(1-p) + 1^2(p) - (0(1-p) + 1(p))^2 \\
    &= p - p^2 \\
    &= p(1-p)
\end{align*}

\paragraph*{Example:}
Let $X$ be a discrete random variable with $\mu = 2$ and variance of $\sigma^2 = 6$. (a) Find $E(X^2)$. (b) Find $V[(X-1)^2]$ and (c) 
Find $\sigma_X$.

\begin{align*}
    (a) \quad E(X^2) &= V(X) + [E(X)]^2 \\
    &= 6 + 2^2 \\
    &= 10 \\
    \\
    (b) \quad V[(X-1)^2] &= E[(X-1)^4] - [E(X-1)^2]^2 \\
    &= E[X^4 - 4X^3 + 6X^2 - 4X + 1] - [E(X^2) - 2E(X) + 1]^2 \\
    &= E[X^4] - 4E[X^3] + 6E[X^2] - 4E[X] + 1 - [10 - 4 + 1]^2 \\
    &= E[X^4] - 4E[X^3] + 60 - 8 + 1 - 49 \\
    &= E[X^4] - 4E[X^3] + 4 \\
    \\
    (c) \quad \sigma_X &= \sqrt{V(X)} = \sqrt{6} \approx 2.45
\end{align*}